\documentclass[a4paper,12pt]{article}

\usepackage[T2A]{fontenc}
\usepackage[utf8]{inputenc}
\usepackage[russian]{babel}

\usepackage{amsmath,amssymb}
\renewcommand{\familydefault}{\sfdefault}
\usepackage{icomma}
\usepackage[a4paper,margin=20mm,top=20mm,bottom=22mm,headheight=25pt,headsep=6mm]{geometry}
\usepackage[nopatch=footnote]{microtype}

\usepackage{tikz}
\usetikzlibrary{calc,arrows.meta,positioning,shapes.geometric}
\usepackage{tkz-euclide}

\usepackage{tabularx,booktabs}
\usepackage{enumitem}
\usepackage{hyperref}
\usepackage{parskip}
\usepackage{fancyhdr}
\usepackage{setspace}
\usepackage{xcolor}
\usepackage{needspace}

% -------------------------------------------------
% Цветовая палитра
% -------------------------------------------------
\definecolor{accent}{HTML}{008080}
\definecolor{darkaccent}{HTML}{004D4D}
\definecolor{textgray}{HTML}{333333}
\definecolor{softgray}{HTML}{6B7280}
\definecolor{lightline}{HTML}{D7E1E1}

\color{textgray}
\setstretch{1.35}
\raggedbottom
\clubpenalty=10000
\widowpenalty=10000
\displaywidowpenalty=10000

\hypersetup{
    colorlinks=true,
    linkcolor=darkaccent,
    urlcolor=darkaccent
}

\setlength{\parindent}{0pt}
\setlist[itemize]{leftmargin=1.6em,itemsep=0.25em,topsep=0.35em}
\setlist[enumerate]{leftmargin=1.8em,itemsep=0.55em,topsep=0.45em}

\pagestyle{fancy}
\fancyhf{}
\fancyhead[L]{\small\color{softgray}Векторы}
\fancyhead[R]{\small\color{softgray}ЕГЭ}
\fancyfoot[C]{\small\color{softgray}\thepage}
\renewcommand{\headrulewidth}{0.3pt}
\renewcommand{\headrule}{%
    \hbox to\headwidth{\color{lightline}\leaders\hrule height \headrulewidth\hfill}%
}

% -------------------------------------------------
% Служебные команды
% -------------------------------------------------
\newcommand{\vect}[1]{\overrightarrow{#1}}
\newcommand{\vct}[1]{\vec{#1}}
\newcommand{\modv}[1]{\left|#1\right|}
\newcommand{\sectionline}{%
    \vspace{0.1em}
    {\color{accent}\rule{\linewidth}{0.6pt}}
    \vspace{0.3em}
}
\newcommand{\topic}[1]{%
    \Needspace{7\baselineskip}%
    \section*{\color{darkaccent}#1}
    \sectionline
}
\newcommand{\checkitem}[1]{%
    \par\noindent
    {\color{accent}\(\square\)}\;#1\par
}

% -------------------------------------------------
% Стили блок-схем
% -------------------------------------------------
\tikzset{
    flowstart/.style={
        ellipse,
        draw=darkaccent,
        line width=0.9pt,
        minimum width=34mm,
        minimum height=10mm,
        align=center,
        font=\small
    },
    flowaction/.style={
        rounded corners=3pt,
        rectangle,
        draw=accent,
        line width=0.9pt,
        text width=82mm,
        minimum height=12mm,
        align=center,
        inner sep=4pt,
        font=\small
    },
    flowbranch/.style={
        rounded corners=3pt,
        rectangle,
        draw=accent,
        line width=0.9pt,
        text width=57mm,
        minimum height=17mm,
        align=center,
        inner sep=4pt,
        font=\small
    },
    flowdecision/.style={
        diamond,
        aspect=2.8,
        draw=darkaccent,
        line width=0.9pt,
        text width=40mm,
        align=center,
        inner sep=2pt,
        font=\small
    },
    flowarrow/.style={
        -{Stealth[length=3mm,width=2mm]},
        line width=0.9pt,
        draw=textgray
    },
    flowlabel/.style={
        font=\small,
        fill=white,
        inner sep=1.5pt
    }
}

% -------------------------------------------------
% Документ
% -------------------------------------------------
\begin{document}

% =================================================
% ТИТУЛЬНЫЙ ЛИСТ
% =================================================
\begin{titlepage}
\thispagestyle{empty}

\begin{center}

\vspace*{25mm}

{\color{darkaccent}\Large\bfseries ЧЕК-ЛИСТ}

\vspace{9mm}

{\Huge\bfseries Векторы}

\vspace{4mm}

{\Large
Координаты, длина, действия,\\
скалярное произведение и угол между векторами
}

\vspace{17mm}

{\color{accent}\rule{0.68\linewidth}{0.9pt}}

\vspace{15mm}

{\large
ЕГЭ по профильной математике
}

\vspace{8mm}

{\large
Формат тренировки: Путь А
}

\vspace{20mm}

{\large
Дата: 23 сентября 2026 г.
}

\vfill

{\large
Лёвин Артём Александрович
}

\vspace{2mm}

{\large
эксклюзивно для Екатерины Гнедковой
}

\vspace{24mm}

\end{center}

\begin{tikzpicture}[remember picture,overlay]
    \draw[
        accent,
        line width=1.1pt
    ]
    ($(current page.south west)+(12mm,14mm)$)
    .. controls
    ($(current page.south west)+(65mm,27mm)$)
    and
    ($(current page.south east)+(-62mm,5mm)$)
    ..
    ($(current page.south east)+(-12mm,16mm)$);
\end{tikzpicture}

\end{titlepage}

% =================================================
% ОСНОВНОЙ ЧЕК-ЛИСТ
% =================================================
\topic{1. Координаты вектора}

Пусть
\[
A(x_A;y_A),
\qquad
B(x_B;y_B).
\]

Тогда координаты вектора \(\vect{AB}\) равны
\[
\boxed{
\vect{AB}
=
(x_B-x_A\,;\,y_B-y_A)
}.
\]

Главное правило:
\[
\boxed{\text{конец} - \text{начало}}
\]
по каждой координате.

\subsection*{\color{darkaccent}Геометрический смысл}

Координаты вектора показывают смещение:

\begin{itemize}
    \item первая координата --- смещение по оси \(x\);
    \item вторая координата --- смещение по оси \(y\).
\end{itemize}

Например, если вектор смещает точку на \(5\) единиц вправо и на \(2\) единицы вверх, то
\[
\vct a=(5;2).
\]

\subsection*{\color{darkaccent}Пример}

Пусть
\[
A(2;2),
\qquad
B(7;4).
\]

Тогда
\[
\vect{AB}
=
(7-2\,;\,4-2)
=
(5;2).
\]

\checkitem{При нахождении координат вектора Вы вычитаете координаты начала из координат конца.}

\checkitem{Отрицательная координата допустима: она означает смещение в отрицательном направлении соответствующей оси.}

% =================================================
\topic{2. Длина вектора}

Если \(\vct a=(x;y)\), то
\[
\boxed{
\modv{\vct a}=\sqrt{x^2+y^2}
}.
\]

Формула следует из теоремы Пифагора. Например,
\[
\vect{AB}=(5;2)
\quad\Longrightarrow\quad
\modv{\vect{AB}}=\sqrt{5^2+2^2}=\sqrt{29}.
\]

Для любого вектора
\[
\modv{\vct a}\ge0,
\qquad
\modv{\vct a}=0\quad\Longleftrightarrow\quad\vct a=\vct 0.
\]

\checkitem{Под корнем складываются квадраты координат.}

\checkitem{Если координата отрицательна, при возведении в квадрат заключайте её в скобки: например, \((-3)^2=9\).}

% =================================================
\topic{3. Умножение вектора на число}

Если
\[
\vct a=(x;y),
\]
то
\[
\boxed{
\lambda\vct a=(\lambda x;\lambda y)
}.
\]

Например,
\[
\vct a=(5;2)
\quad\Longrightarrow\quad
5\vct a=(25;10).
\]

При умножении на \(-1\) получаем противоположный вектор:
\[
-\vct a=(-x;-y),
\qquad
\modv{-\vct a}=\modv{\vct a}.
\]
Для ненулевого \(\vct a\) векторы \(\vct a\) и \(-\vct a\) имеют противоположные направления.

\checkitem{Число умножает каждую координату вектора.}

% =================================================
\topic{4. Сложение, вычитание и линейная комбинация}

Пусть
\[
\vct a=(x_a;y_a),
\qquad
\vct b=(x_b;y_b).
\]

Тогда
\[
\boxed{
\begin{aligned}
\vct a+\vct b&=(x_a+x_b\,;\,y_a+y_b),\\
\vct a-\vct b&=(x_a-x_b\,;\,y_a-y_b).
\end{aligned}
}
\]

Также
\[
\vct a-\vct b=\vct a+(-\vct b).
\]

\subsection*{\color{darkaccent}Пример: линейная комбинация}

Пусть \(\vct k=(1;2)\), \(\vct t=(3;4)\). Найдём длину вектора \(3\vct k-2\vct t\):
\[
\begin{aligned}
3\vct k&=(3;6),
&2\vct t&=(6;8),\\
3\vct k-2\vct t&=(-3;-2),
&\modv{3\vct k-2\vct t}&=\sqrt{13}.
\end{aligned}
\]

\subsection*{\color{darkaccent}Надёжный порядок действий}

\[
\boxed{
\text{коэффициенты}
\;\longrightarrow\;
\text{координаты результата}
\;\longrightarrow\;
\text{длина при необходимости}
}
\]

Такой порядок позволяет отдельно контролировать каждый этап вычислений.

% =================================================
% ОТДЕЛЬНАЯ СТРАНИЦА БЛОК-СХЕМЫ
% =================================================
\clearpage
\thispagestyle{fancy}

\begin{center}
{\Large\bfseries\color{darkaccent}
Алгоритм: линейная комбинация векторов
}
\end{center}

\vfill

\begin{center}
\begin{tikzpicture}[node distance=8mm]

\node[flowstart] (start) {Начало};

\node[flowaction,below=of start] (write)
{Выпишите координаты всех исходных векторов};

\node[flowaction,below=of write] (multiply)
{Умножьте обе координаты каждого вектора на его числовой коэффициент};

\node[flowaction,below=of multiply] (combine)
{Сложите или вычтите соответствующие координаты};

\node[flowdecision,below=11mm of combine] (target)
{Требуется длина полученного вектора?};

\node[flowbranch,below=16mm of target,xshift=-40mm] (length)
{Да: для \(\vct c=(x;y)\) вычислите
\(\modv{\vct c}=\sqrt{x^2+y^2}\)};

\node[flowbranch,below=16mm of target,xshift=40mm] (coords)
{Нет: оставьте найденные координаты вектора как окончательный результат};

\coordinate (merge) at ($(length.south)!0.5!(coords.south)+(0,-13mm)$);
\node[flowstart,below=0mm of merge] (finish) {Ответ};

\draw[flowarrow] (start) -- (write);
\draw[flowarrow] (write) -- (multiply);
\draw[flowarrow] (multiply) -- (combine);
\draw[flowarrow] (combine) -- (target);

\draw[flowarrow]
(target.south west)
|- node[flowlabel,pos=0.3,left]{да}
(length.north);

\draw[flowarrow]
(target.south east)
|- node[flowlabel,pos=0.3,right]{нет}
(coords.north);

\draw[flowarrow] (length.south) |- (finish.west);
\draw[flowarrow] (coords.south) |- (finish.east);

\end{tikzpicture}
\end{center}

\vfill
\clearpage

% =================================================
\topic{5. Скалярное произведение по координатам}

Пусть
\[
\vct a=(x_a;y_a),
\qquad
\vct b=(x_b;y_b).
\]

Тогда
\[
\boxed{
\vct a\cdot\vct b
=
x_ax_b+y_ay_b
}.
\]

Результат скалярного произведения --- \textbf{число}.

\subsection*{\color{darkaccent}Пример}

Пусть
\[
\vct a=(2;-3),
\qquad
\vct b=(1;-5).
\]

Тогда
\[
\vct a\cdot\vct b
=
2\cdot1+(-3)\cdot(-5)
=
2+15
=
17.
\]

\checkitem{Перемножаются соответствующие координаты, после чего два произведения складываются.}

\checkitem{Особенно внимательно проверяйте знак произведения двух отрицательных чисел.}

% =================================================
\topic{6. Скалярное произведение и угол между векторами}

Если \(\varphi\) --- угол между ненулевыми векторами \(\vct a\) и \(\vct b\), где
\[
0^\circ\le\varphi\le180^\circ,
\]
то
\[
\boxed{
\vct a\cdot\vct b
=
\modv{\vct a}\,\modv{\vct b}\cos\varphi
}.
\]

Поскольку оба вектора ненулевые,
\[
\modv{\vct a}>0,
\qquad
\modv{\vct b}>0,
\]
поэтому можно разделить на произведение их длин:
\[
\boxed{
\cos\varphi
=
\frac{\vct a\cdot\vct b}
{\modv{\vct a}\modv{\vct b}}
}.
\]

\subsection*{\color{darkaccent}Если известны координаты}

Удобный порядок вычислений:
\[
\vct a\cdot\vct b
\quad\longrightarrow\quad
\modv{\vct a}
\quad\longrightarrow\quad
\modv{\vct b}
\quad\longrightarrow\quad
\cos\varphi.
\]

\subsection*{\color{darkaccent}Пример}

Пусть
\[
\vct a=(3;4),
\qquad
\vct b=(4;-3).
\]

Тогда
\[
\vct a\cdot\vct b
=
3\cdot4+4\cdot(-3)
=
0.
\]

При этом
\[
\modv{\vct a}=5,
\qquad
\modv{\vct b}=5.
\]

Следовательно,
\[
\cos\varphi
=
\frac{0}{5\cdot5}
=
0,
\]
значит,
\[
\varphi=90^\circ.
\]

Для ненулевых векторов получаем важный признак:
\[
\boxed{
\vct a\cdot\vct b=0
\quad\Longleftrightarrow\quad
\vct a\perp\vct b
}.
\]

Кроме того, знак скалярного произведения помогает быстро проверить характер угла:
\[
\begin{aligned}
\vct a\cdot\vct b&>0 &&\Longrightarrow && 0^\circ\le\varphi<90^\circ,\\
\vct a\cdot\vct b&=0 &&\Longrightarrow && \varphi=90^\circ,\\
\vct a\cdot\vct b&<0 &&\Longrightarrow && 90^\circ<\varphi\le180^\circ.
\end{aligned}
\]

% =================================================
\topic{7. Как выбрать формулу скалярного произведения}

Ориентируйтесь на данные задачи.

Если известны координаты обоих векторов, удобно начать с
\[
\vct a\cdot\vct b=x_ax_b+y_ay_b.
\]

Если известны длины векторов и угол между ними, используйте
\[
\vct a\cdot\vct b
=
\modv{\vct a}\modv{\vct b}\cos\varphi.
\]

Если известны координаты ненулевых векторов, а требуется угол, формулы объединяются:
\[
\boxed{
\cos\varphi
=
\frac{x_ax_b+y_ay_b}
{\sqrt{x_a^2+y_a^2}\,\sqrt{x_b^2+y_b^2}}
}.
\]

\subsection*{\color{darkaccent}Контроль результата}

Для угла между векторами обязательно
\[
-1\le\cos\varphi\le1.
\]

Если после вычислений получено число вне этого диапазона, в вычислениях есть ошибка.

% =================================================
\topic{8. Сложение векторов геометрически}

Для сложения векторов действует \textbf{правило треугольника}: если конец первого вектора совместить с началом второго, то вектор суммы направлен от начала первого вектора к концу второго.

В частности, для точек \(A\), \(B\), \(C\)
\[
\boxed{
\vect{AB}+\vect{BC}=\vect{AC}
}.
\]

При вычитании удобно использовать равенство
\[
\vct a-\vct b=\vct a+(-\vct b)
\]
и затем применять правило треугольника.

Для ненулевого \(\vct b\) векторы \(\vct b\) и \(-\vct b\) имеют одинаковую длину и противоположные направления.

% =================================================
\topic{9. Геометрическая задача с ромбом}

\begin{minipage}[t]{0.59\linewidth}
Пусть \(ABCD\) --- ромб, его диагонали \(AC\) и \(BD\) пересекаются в точке \(O\). Диагонали ромба точкой пересечения делятся пополам и взаимно перпендикулярны.

Рассмотрим случай
\[
AC=12,
\qquad
BD=16.
\]
Тогда
\[
AO=OC=6,
\qquad
BO=OD=8.
\]
\end{minipage}
\hfill
\begin{minipage}[t]{0.36\linewidth}
\centering
\vspace{0pt}
\begin{tikzpicture}[scale=0.47]
\tkzSetUpLine[line width=1pt,color=accent]
\tkzSetUpPoint[color=accent,fill=accent,size=3]
\tkzSetUpLabel[color=black,font=\small]

\tkzDefPoint(0,4){A}
\tkzDefPoint(-5,0){B}
\tkzDefPoint(0,-4){C}
\tkzDefPoint(5,0){D}
\tkzDefPoint(0,0){O}

\tkzDrawPolygon(A,B,C,D)
\tkzDrawSegments(A,C B,D A,D)
\tkzDrawPoints(A,B,C,D,O)

\tkzLabelPoints[above](A)
\tkzLabelPoints[left](B)
\tkzLabelPoints[below](C)
\tkzLabelPoints[right](D)
\tkzLabelPoints[below right](O)

\tkzMarkRightAngle[size=0.38](A,O,D)
\end{tikzpicture}
\end{minipage}

Поскольку \(B\), \(O\), \(D\) лежат на одной прямой в порядке \(B\! - \!O\! - \!D\) и \(BO=OD\),
\[
\vect{BO}=\vect{OD}.
\]
По правилу треугольника
\[
\vect{AO}+\vect{BO}
=
\vect{AO}+\vect{OD}
=
\vect{AD}.
\]
Следовательно,
\[
\modv{\vect{AO}+\vect{BO}}
=AD
=\sqrt{AO^2+OD^2}
=\sqrt{6^2+8^2}
=10.
\]

\checkitem{Если координатная система в геометрической задаче не задана, её можно ввести самостоятельно, когда это упрощает решение.}

\checkitem{Если в качестве осей выбираются диагонали ромба, их перпендикулярность гарантирует корректную прямоугольную систему координат.}

% =================================================
\topic{10. Квадрат длины вектора}

Если требуется найти квадрат длины вектора, извлекать корень не нужно.

Для
\[
\vct a=(x;y)
\]
имеем
\[
\boxed{
\modv{\vct a}^{2}=x^2+y^2
}.
\]

Например, если
\[
\vct a=(10;10),
\]
то
\[
\modv{\vct a}^{2}
=
10^2+10^2
=
200.
\]

Это удобно в задачах, где требуется именно \(\modv{\vct a}^{2}\): лишнее извлечение корня увеличивает объём вычислений.

% =================================================
\topic{11. Типичные ошибки}

\begin{enumerate}
    \item \textbf{Перепутан порядок при поиске координат.}

    Для \(\vect{AB}\) нужно считать
    \[
    B-A,
    \]
    то есть координаты конца минус координаты начала.

    \item \textbf{Числовой множитель применён только к одной координате.}

    Если
    \[
    \vct a=(x;y),
    \]
    то
    \[
    3\vct a=(3x;3y).
    \]

    \item \textbf{При вычитании векторов потерян знак.}

    Например,
    \[
    (3;6)-(6;8)=(-3;-2).
    \]

    \item \textbf{Отрицательная координата неверно возведена в квадрат.}

    Верно:
    \[
    (-3)^2=9.
    \]

    \item \textbf{Смешаны скалярное произведение и длина.}

    Выражение
    \[
    \vct a\cdot\vct b
    \]
    является числом, а
    \[
    \modv{\vct a}
    \]
    --- длиной одного вектора.

    \item \textbf{Формула через угол используется при нулевом векторе.}

    Формула
    \[
    \cos\varphi
    =
    \frac{\vct a\cdot\vct b}
    {\modv{\vct a}\modv{\vct b}}
    \]
    применима только к двум ненулевым векторам.

    \item \textbf{Сложное выражение считается одним действием.}

    Надёжнее отдельно получить, например,
    \[
    3\vct a,
    \qquad
    2\vct b,
    \]
    затем найти их разность и только после этого вычислять длину.
\end{enumerate}

% =================================================
\topic{12. Итоговый чек-лист}

Перед тренировкой убедитесь, что Вы умеете:

\checkitem{находить координаты вектора по координатам его начала и конца;}

\checkitem{понимать координаты вектора как смещение по двум осям;}

\checkitem{вычислять длину вектора и квадрат его длины;}

\checkitem{умножать вектор на число;}

\checkitem{складывать и вычитать векторы по координатам;}

\checkitem{последовательно вычислять линейную комбинацию вида \(m\vct a+n\vct b\);}

\checkitem{находить скалярное произведение по координатам;}

\checkitem{использовать формулу скалярного произведения через длины и угол;}

\checkitem{находить \(\cos\varphi\) между двумя ненулевыми векторами;}

\checkitem{распознавать перпендикулярность ненулевых векторов по условию \(\vct a\cdot\vct b=0\);}

\checkitem{складывать векторы геометрически по правилу треугольника;}

\checkitem{заменять вычитание сложением с противоположным вектором;}

\checkitem{вводить удобную систему координат в геометрической задаче.}

% =================================================
% ТРЕНИРОВОЧНЫЙ БЛОК
% =================================================
\clearpage

\topic{Тренировочный блок --- Путь А}

\textbf{7 заданий.} Решайте последовательно и записывайте промежуточные координаты векторов.

\begin{enumerate}

    \item Даны точки
    \[
    A(-3;2),
    \qquad
    B(4;-2).
    \]
    Найдите координаты вектора \(\vect{AB}\) и его длину.

    \vspace{8mm}

    \item Даны векторы
    \[
    \vct a=(2;-1),
    \qquad
    \vct b=(-3;4).
    \]
    Найдите координаты и длину вектора
    \[
    \vct c=3\vct a-2\vct b.
    \]

    \vspace{8mm}

    \item Даны векторы
    \[
    \vct a=(4;-2),
    \qquad
    \vct b=(1;5).
    \]
    Найдите
    \[
    \vct a\cdot\vct b.
    \]

    \vspace{8mm}

    \item Даны векторы
    \[
    \vct a=(3;4),
    \qquad
    \vct b=(4;-3).
    \]
    Найдите косинус угла между ними и величину этого угла.

    \vspace{8mm}

    \item Даны точки
    \[
    A(1;-2),\quad B(5;1),
    \qquad
    C(-1;3),\quad D(2;-1).
    \]
    Найдите
    \[
    \modv{\vect{AB}+\vect{CD}}^2.
    \]

    \vspace{8mm}

    \item Известно, что
    \[
    \modv{\vct a}=6,
    \qquad
    \modv{\vct b}=4,
    \]
    а угол между векторами равен
    \[
    60^\circ.
    \]
    Найдите
    \[
    \vct a\cdot\vct b.
    \]

    \vspace{8mm}

    \item Диагонали ромба \(ABCD\) пересекаются в точке \(O\). Известно, что
    \[
    AC=10,
    \qquad
    BD=24.
    \]
    Найдите
    \[
    \modv{\vect{AO}+\vect{BO}}.
    \]

\end{enumerate}

% =================================================
% ОТВЕТЫ
% =================================================
\clearpage

\topic{Ответы}

\renewcommand{\arraystretch}{1.45}

\begin{center}
\begin{tabularx}{\linewidth}{
    >{\centering\arraybackslash}p{14mm}
    X
}
\toprule
\textbf{№} & \textbf{Ответ} \\
\midrule
1 & \(\vect{AB}=(7;-4),\quad \modv{\vect{AB}}=\sqrt{65}\) \\
2 & \(\vct c=(12;-11),\quad \modv{\vct c}=\sqrt{265}\) \\
3 & \(-6\) \\
4 & \(\cos\varphi=0,\quad \varphi=90^\circ\) \\
5 & \(50\) \\
6 & \(12\) \\
7 & \(13\) \\
\bottomrule
\end{tabularx}
\end{center}

\vfill

\begin{center}
{\small\color{softgray}
Лёвин Артём Александрович
\quad\textbullet\quad
эксклюзивно для Екатерины Гнедковой
}
\end{center}

\end{document}